% =====================================================================
% Identidade visual de gráficos, tabelas, quadros e guias de escrita
% Adaptado do template ufal-tesis (OA-RAG / Hossain e Sarıyüce)
% =====================================================================

\usepackage{colortbl}
\usepackage{pifont}
\usepackage{fontawesome5}   % ícones para diagramas e fluxogramas
\usepackage{fvextra}        % listagem de prompts/arquivos nos apêndices
\usepackage{placeins}       % \FloatBarrier: mantém figuras dentro da própria seção
\tcbuselibrary{breakable}
\usetikzlibrary{patterns, shapes.arrows}
\usepgfplotslibrary{groupplots, statistics}

% Quadros numerados por capítulo (Quadro 4.1, 5.1...): \begin{quadro}[label=...]{Título}
\newtcolorbox[auto counter, number within=chapter]{quadro}[2][]{
  colback=gray!5,
  colframe=gray!60,
  fonttitle=\bfseries,
  title={Quadro~\thetcbcounter: #2},
  #1
}

% --- Guias de escrita do esqueleto ------------------------------------
% Para ocultar todos os quadros de roteiro na compilação final, altere
% \guiaescritatrue para \guiaescritafalse.
\newif\ifguiaescrita
\guiaescritafalse
\newcommand{\guiasecao}[1]{%
  \ifguiaescrita
    \begin{tcolorbox}[
      breakable,
      colback=vAzul!4,
      colframe=vAzul!45,
      boxrule=0.5pt,
      arc=2pt,
      left=6pt, right=6pt, top=4pt, bottom=4pt,
      fonttitle=\sffamily\bfseries\footnotesize,
      coltitle=vAzul!80!black,
      colbacktitle=vAzul!10,
      title={Roteiro da Seção (esqueleto a redigir)}
    ]
      \footnotesize\sffamily #1
    \end{tcolorbox}%
  \fi
}

% --- Paleta principal -------------------------------------------------
\definecolor{vAzul}{HTML}{2A78D6}
\definecolor{vVerde}{HTML}{1BAF7A}
\definecolor{vLaranja}{HTML}{EB6834}
\definecolor{vAmbar}{HTML}{EDA100}
\definecolor{vRoxo}{HTML}{4A3AA7}
\definecolor{vVermelho}{HTML}{E34948}
\definecolor{vTeal}{HTML}{2A9D8F}
\definecolor{vCinza}{HTML}{9B9A94}
\definecolor{vNeutro}{HTML}{DDDCD5}

% Cores semânticas sugeridas para vereditos do Judge
\colorlet{corAC}{vVerde}
\colorlet{corWA}{vVermelho}
\colorlet{corCE}{vRoxo}
\colorlet{corRE}{vLaranja}
\colorlet{corTLE}{vAmbar}
\colorlet{corMLE}{vTeal}

% --- Tabelas ----------------------------------------------------------
\colorlet{linhaDestaque}{black!6}
\colorlet{colunaDestaque}{black!8}
\newcommand{\tfirst}[1]{{\color{red!80!black}#1}}
\newcommand{\tsecond}[1]{{\color{blue!80!black}#1}}
\newcommand{\gain}[1]{\;{\color{green!50!black}\ensuremath{\uparrow}\,#1}}
\newcommand{\loss}[1]{\;{\color{red!75!black}\ensuremath{\downarrow}\,#1}}
\newcommand{\cmark}{\ding{51}}
\newcommand{\nomark}{{\color{black!35}--}}
\newcommand{\tcok}{\textcolor{green!50!black}{\ding{51}}}
\newcommand{\tcno}{\textcolor{red!75!black}{\ding{55}}}

% Amostra de cor para legendas
\newcommand{\chip}[1]{\tikz[baseline=-0.6ex]\filldraw[fill=#1, draw=black!40, line width=0.3pt] (0,-0.09) rectangle (0.30,0.09);\hspace{2.5pt}}

% Barra proporcional dentro de célula de tabela (valor em [0,1]):
% \barcell{cor}{valor}{rótulo}
\newcommand{\barcell}[3]{%
  \tikz[baseline=-0.6ex]{%
    \fill[black!5] (0,-0.11) rectangle (1.6,0.11);%
    \fill[#1] (0,-0.11) rectangle ({1.6*#2},0.11);%
    \node[anchor=west, inner sep=0pt, font=\footnotesize] at ({1.6*#2+0.06},0) {#3};%
  }%
}

% --- Posicionamento de figuras ----------------------------------------
\renewcommand{\topfraction}{.92}\renewcommand{\dbltopfraction}{.92}
\renewcommand{\textfraction}{.06}\renewcommand{\floatpagefraction}{.72}

% --- Heatmap de células -----------------------------------------------
% Branco/verde (0) -> amarelo (50) -> vermelho (100); imprime o valor.
\newcommand{\heatmap}[1]{%
    \pgfmathsetmacro{\val}{#1}%
    \pgfmathsetmacro{\r}{ifthenelse(\val<50, \val/50, 1)}%
    \pgfmathsetmacro{\g}{ifthenelse(\val<50, 1, 2-\val/50)}%
    \pgfmathsetmacro{\b}{0}%
    \edef\temp{\noexpand\cellcolor[rgb]{\r,\g,\b}}%
    \temp #1%
}

% --- Número do capítulo em fonte monoespaçada no cabeçalho de abertura -
\usepackage{etoolbox}
\makeatletter
\patchcmd{\@makechapterhead}{\thechapter}{{\ttfamily\fontsize{34}{34}\selectfont\thechapter}}{}{\PackageWarning{estilo-visual}{patch do numero do capitulo falhou}}
\makeatother
